% =====================================================================
% Documento de previsualización de los borradores (NO es la memoria).
% Sirve solo para compilar los capítulos sueltos mientras no exista la
% plantilla oficial de la ESI (P-04). Al integrar la plantilla, copiar a su
% preámbulo los paquetes que falten y hacer \input de cada capítulo.
%   latexmk -pdf borradores.tex
% =====================================================================
\documentclass[11pt,a4paper,openany]{report}

\usepackage[utf8]{inputenc}
\usepackage[T1]{fontenc}
\usepackage[spanish,es-tabla]{babel}
\usepackage{lmodern}
\usepackage[margin=2.5cm]{geometry}
\usepackage{booktabs}
\usepackage{tabularx}
\usepackage{array}
\usepackage{enumitem}
\usepackage{xcolor}
\usepackage{tikz}
\usetikzlibrary{positioning,arrows.meta,shapes.geometric,fit,backgrounds,calc}
\usepackage[colorinlistoftodos,spanish]{todonotes}
\setlength{\marginparwidth}{2cm}
\usepackage{hyperref}

\title{PymeKit --- Borradores de la memoria}
\date{\today}

\begin{document}
\maketitle
\listoftodos[Pendientes del autor]

\setcounter{chapter}{1}
% =====================================================================
% Capítulo 2 · Planificación
% Borrador generado a partir de docs/tfg/{PLAN,PROGRESO,BITACORA,DECISIONES}.md,
% AGENTS.md, .claude/, .github/workflows/ci.yml y el historial de git.
% Estado de las fuentes: 2026-09-29 (F2.4b cerrada).
% Paquetes necesarios: booktabs, tabularx, todonotes, hyperref.
% =====================================================================

\chapter{Planificación}
\label{cap:planificacion}

En este capítulo se describe cómo se ha organizado el trabajo: la metodología
seguida, la descomposición en fases y su relación con la Estructura de
Descomposición del Trabajo (EDT) de la propuesta, la temporización, la gestión
de riesgos, las herramientas empleadas y la estructura de costes.

% ---------------------------------------------------------------------
\section{Metodología}
\label{sec:metodologia}

\subsection{Desarrollo incremental por fases}

El proyecto se ha abordado con un modelo de desarrollo \textbf{incremental}:
el trabajo se divide en fases (F0--F8) y, dentro de las fases más grandes, en
incrementos más pequeños (por ejemplo, F2.1--F2.9 en la integración del CMS).
Cada fase o incremento tiene tres elementos fijados de antemano en el plan del
proyecto:

\begin{itemize}
  \item un \textbf{objetivo} y una lista de entregables;
  \item un \textbf{criterio de «hecho»} verificable (por ejemplo, «la
        aplicación arranca en local, CU-01 funciona a mano y la integración
        continua está en verde»);
  \item las \textbf{dependencias} con otras fases.
\end{itemize}

Un incremento solo se da por cerrado cuando supera su criterio de «hecho» y la
verificación común descrita en la sección~\ref{sec:verificacion-proceso}. Así,
al final de cada incremento el sistema queda en un estado funcional y
comprobado, y los retrasos afectan al alcance de los incrementos posteriores,
no a la estabilidad de lo ya construido. Esta forma de trabajar encaja con el
objetivo del TFG, una arquitectura \emph{reutilizable}: cada incremento añade
un módulo que debe poder integrarse sin romper los anteriores (RNF-04,
RNF-06).

Las decisiones de arquitectura o de alcance se registran en el momento en que
se toman como \textit{Architecture Decision Records} (ADR, registros de
decisiones de arquitectura)~\cite{nygard2011adr}, con su contexto, las
alternativas consideradas y sus consecuencias. En el momento de redactar este
borrador hay dieciséis (ADR-001 a ADR-016); los relevantes para el diseño se
resumen en el capítulo~\ref{cap:diseno}. Las incidencias (fallos propios o
heredados, problemas de entorno y de herramientas) se anotan también en el
momento en una bitácora con el formato \emph{qué pasó, causa, solución,
evidencia y lección}; esa bitácora es la fuente de los riesgos materializados
de la sección~\ref{sec:riesgos}.

\subsection{Desarrollo asistido por agentes de IA}
\label{sec:metodologia-ia}

% PENDIENTE (autor): comprobar si la ESI/UCA exige una declaración formal del
% uso de herramientas de IA generativa en el TFG (formato y ubicación) y
% ajustar este apartado a esa normativa.

Una parte sustancial del código, de las pruebas y de la documentación se ha
producido con la ayuda de un asistente de programación basado en modelos de
lenguaje que actúa como \textit{agente} (un programa que, a partir de una
instrucción, lee ficheros, ejecuta órdenes y propone o aplica cambios en el
repositorio). En concreto se ha usado Claude Code~\cite{claudecode}. Todos los
\textit{commits} del repositorio registran esa coautoría mediante la línea
\texttt{Co-Authored-By} del mensaje.
\todo{Confirmar el modelo o modelos de IA empleados y el periodo de uso para
la declaración de uso de IA.}

Para que este modo de trabajo sea controlable y evaluable, antes de escribir
código de la aplicación (fase F0) se construyó un \textit{harness}, es decir,
un arnés de directrices y controles automáticos que acota lo que hacen los
agentes y comprueba el resultado. Sus componentes son:

\begin{itemize}
  \item \textbf{Directrices escritas.} Un fichero raíz de directrices para
        agentes y desarrolladores (\texttt{AGENTS.md}) y otros por paquete que
        fijan el \textit{stack}, las convenciones, las reglas de seguridad de
        la base de datos, el idioma de comentarios y documentación (ADR-003) y
        una lista de verificación obligatoria al terminar cualquier cambio.
  \item \textbf{Documentación de seguimiento} en \texttt{docs/tfg/}: plan,
        requisitos, decisiones, trazabilidad, progreso, dedicación y bitácora.
        Un agente debe leer el progreso y el plan de la fase antes de empezar,
        y no puede reabrir una decisión sin registrar un ADR nuevo.
  \item \textbf{Procedimientos reutilizables} (\textit{skills}) para tareas
        repetitivas o delicadas: portar un módulo, reescribir comentarios,
        registrar una decisión, revisar un cambio de forma adversarial, auditar
        las políticas de seguridad de la base de datos o redactar un borrador
        de la memoria.
  \item \textbf{Controles automáticos} (\textit{hooks}) que se ejecutan sin
        intervención: al inicio de cada sesión se inyecta la fase actual y sus
        tareas pendientes; tras cada edición se comprueba que el fichero no
        contiene referencias prohibidas (ADR-007); y antes de terminar se
        recuerda actualizar el progreso si hay cambios de código sin
        registrar.
  \item \textbf{Permisos} explícitos: el código de partida es de solo lectura
        para los agentes y no pueden leer ficheros de secretos.
\end{itemize}

\paragraph{Reparto de responsabilidades.}
El reparto entre el autor y los agentes ha sido el siguiente:

\begin{itemize}
  \item \textbf{El autor} fija el alcance y los requisitos, toma las decisiones
        de arquitectura (por ejemplo, la integración del CMS en la consola de
        administración, ADR-011, fue una propuesta suya), revisa y acepta o
        rechaza los cambios, realiza pruebas manuales exploratorias y es el
        único que autoriza los \textit{commits} y su publicación. Las
        directrices prohíben a los agentes hacer \textit{commit} o
        \textit{push} sin su orden expresa.
  \item \textbf{Los agentes} analizan el código existente, proponen
        alternativas, implementan los incrementos, escriben las pruebas,
        ejecutan la verificación, redactan la documentación de seguimiento y
        los borradores de la memoria, y realizan las revisiones adversariales.
\end{itemize}
\todo{El autor debe revisar y completar este reparto con su propia valoración:
qué partes ha escrito o reescrito a mano, qué criterio ha seguido para aceptar
cambios y cuánto tiempo ha dedicado a la revisión.}

\paragraph{Cómo se verifica el trabajo.}
\label{sec:verificacion-proceso}
Ningún cambio se acepta por el mero hecho de haberlo producido un agente. Todo
cambio pasa por la misma cadena de comprobaciones, que el
\texttt{AGENTS.md} declara obligatoria:

\begin{enumerate}
  \item comprobación de tipos (TypeScript estricto), \textit{lint} y formato;
  \item comprobación de desmarcado (\texttt{check-branding});
  \item pruebas afectadas: unitarias, de base de datos (pgTAP) y
        \textit{end-to-end} (E2E) si se toca un flujo de usuario;
  \item revisión adversarial del \textit{diff} en dos pistas paralelas
        (calidad y arquitectura; búsqueda de errores y fallos de seguridad);
  \item si el cambio afecta a la base de datos, auditoría adversarial de las
        políticas de seguridad a nivel de fila con pruebas de ataque ejecutables
        y una etapa de refutación independiente (véase la
        sección~\ref{sec:verificacion-seguridad});
  \item revisión de los comentarios del código;
  \item actualización del progreso, la trazabilidad y, si procede, las
        decisiones y la bitácora.
\end{enumerate}

Las comprobaciones automatizables se repiten en la integración continua
(sección~\ref{sec:herramientas}). La revisión humana sigue siendo necesaria:
la incidencia B-14 (la consola de administración se cargaba sin segundo
factor) la detectó el autor con una prueba manual que ninguna prueba
automática cubría. De igual modo, la incidencia B-11 muestra que un test en
verde no garantiza la corrección: un test heredado afirmaba precisamente el
comportamiento inseguro. Por eso la metodología combina la verificación
automática con la lectura crítica de las aserciones y con pruebas
exploratorias del autor.

% ---------------------------------------------------------------------
\section{Fases y relación con la EDT}
\label{sec:fases}

La propuesta del TFG define una EDT con seis paquetes de trabajo:
(1)~Análisis, (2)~Arquitectura, (3)~Modelo de datos, (4)~Implementación,
(5)~Validación y (6)~Documentación. La tabla~\ref{tab:fases} relaciona cada
fase con esos paquetes.

\begin{table}[htbp]
  \centering
  \small
  \caption{Fases del proyecto, paquetes de la EDT que cubren y criterio de «hecho».}
  \label{tab:fases}
  \begin{tabularx}{\textwidth}{@{}l>{\raggedright\arraybackslash}p{3.4cm}cX@{}}
    \toprule
    \textbf{Fase} & \textbf{Nombre} & \textbf{EDT} & \textbf{Hecho cuando\ldots} \\
    \midrule
    F0 & Harness y planificación & 1 & El control de desmarcado pasa, los \textit{hooks} funcionan y existe el primer \textit{commit} \\
    F1 & Base SaaS importada y funcionando & 4 & La aplicación arranca en local, CU-01 funciona a mano y la CI está en verde \\
    F2 & Integración del CMS & 2, 3, 4 & Cada incremento F2.1--F2.9 cumple su criterio (BD, API, interfaz, explorador, usuarios y almacenamiento, auditoría, RBAC, paneles, cierre) \\
    F3 & Desmarcado completo e i18n en español & 4 & Control de desmarcado sin excepciones e interfaz en español por defecto \\
    F4 & Comentarios didácticos en español & 4, 6 & Sin comentarios en inglés en los módulos críticos y revisión aprobada \\
    F5 & Adaptación a pymes y configurabilidad & 2, 3, 4 & CU-02 a CU-05 funcionan y la configuración está documentada \\
    F6 & Validación & 5 & CI en verde y resultados documentados \\
    F7 & Despliegue & 4, 6 & Despliegue de demostración accesible y manual probado desde cero \\
    F8 & Memoria y manuales & 1, 2, 3, 6 & Continua; se cierra tras F6 y F7 \\
    \bottomrule
  \end{tabularx}
\end{table}

Las dependencias principales son lineales (F0 $\rightarrow$ F1 $\rightarrow$
F2 $\rightarrow$ F3 $\rightarrow$ F5 $\rightarrow$ F6), con dos excepciones:
F4 puede solaparse con F3 y F5, y F8 se desarrolla de forma continua, de modo
que cada fase cerrada deja un borrador de su parte de la memoria.

La fase F2 es la de mayor esfuerzo. Al comenzarla se evaluaron tres formas de
integrar el CMS (aplicaciones separadas, integración híbrida e integración
total) y se eligió la integración total en la consola de administración
(ADR-011), a pesar de su mayor coste estimado, por dar lugar a un único
servicio, un único inicio de sesión y un único estilo de código. Para
controlar ese riesgo, la fase se dividió en nueve incrementos y se identificó
de antemano el candidato a recortar si el plazo lo exige: los paneles
configurables (RF-11, requisito deseable).

\paragraph{Estado en la fecha de este borrador.}
F0 y F1 están cerradas. En F2 están cerrados los incrementos F2.1 (base de
datos), F2.2 (API), F2.3 (base de la interfaz), F2.4a (listado del explorador
de datos) y F2.4b (ficha de registro).
% PENDIENTE: actualizar este párrafo al integrar el capítulo.

% ---------------------------------------------------------------------
\section{Temporización}
\label{sec:temporizacion}

La dedicación se registra en el documento de progreso del proyecto. Las cifras
de la tabla~\ref{tab:dedicacion} son \textbf{estimaciones} obtenidas a partir
de la hora de los \textit{commits} (desde el inicio de la sesión hasta el
último \textit{commit} de cada tramo), redondeadas al cuarto de hora; el autor
debe revisarlas y registrar las siguientes sesiones con su dedicación real.

\begin{table}[htbp]
  \centering
  \small
  \caption{Dedicación registrada (estimada) por tramos de trabajo.}
  \label{tab:dedicacion}
  \begin{tabularx}{\textwidth}{@{}llXr@{}}
    \toprule
    \textbf{Fecha} & \textbf{Tramo} & \textbf{Fases} & \textbf{Horas} \\
    \midrule
    29/09/2026 & 11:15--11:58 & F0 (harness y plan) & $\approx$ 0,75 \\
    29/09/2026 & 11:58--12:35 & F1 (base SaaS, Docker, CI) & $\approx$ 0,75 \\
    29/09/2026 & 12:35--14:19 & F2 (análisis del CMS y decisión de integración) + F2.1 (BD y auditoría de seguridad) & $\approx$ 1,75 \\
    29/09/2026 & 14:19--15:05 & F2.2 (API del CMS) & $\approx$ 0,75 \\
    29/09/2026 & 15:05--15:50 & F2.3 (base de la interfaz) & $\approx$ 0,75 \\
    29/09/2026 & 15:50--17:02 & Corrección del MFA en la consola + F2.4a (listado) & $\approx$ 1,25 \\
    29/09/2026 & 17:02--17:50 & F2.4b (ficha de registro) + bitácora y trazabilidad & $\approx$ 0,75 \\
    \midrule
    \multicolumn{3}{@{}l}{\textbf{Total registrado}} & $\approx$ \textbf{6,75} \\
    \bottomrule
  \end{tabularx}
\end{table}

\todo{Completar la tabla con las sesiones posteriores y con la dedicación
real revisada por el autor (incluido el tiempo de revisión, pruebas manuales,
reuniones con el tutor y redacción que no deja \textit{commits}).}

\todo{Añadir la planificación inicial estimada por fase (horas previstas) para
compararla con la real, y un diagrama de Gantt (p.\,ej.\ con el paquete
\texttt{pgfgantt}) con las fechas previstas y reales de cada fase.}

% PENDIENTE (autor): las horas registradas corresponden a sesiones de trabajo
% asistido por agentes (sección~\ref{sec:metodologia-ia}). Conviene explicar en
% la memoria qué incluyen y qué no, para que la comparación con la carga de
% créditos del TFG sea honesta.

% ---------------------------------------------------------------------
\section{Gestión de riesgos}
\label{sec:riesgos}

La tabla~\ref{tab:riesgos} recoge los riesgos identificados. Para cada uno se
indica la medida de mitigación y si se ha llegado a materializar, con la
incidencia de la bitácora que lo documenta (identificadores B-xx).

\begin{table}[htbp]
  \centering
  \footnotesize
  \caption{Riesgos del proyecto, mitigación y materialización.}
  \label{tab:riesgos}
  \begin{tabularx}{\textwidth}{@{}l>{\raggedright\arraybackslash}p{3.6cm}>{\raggedright\arraybackslash}Xl@{}}
    \toprule
    \textbf{ID} & \textbf{Riesgo} & \textbf{Mitigación} & \textbf{Materializado} \\
    \midrule
    R-01 & Desbordamiento del plazo por el esfuerzo de integración del CMS &
      Incrementos F2.1--F2.9 funcionales; RF-11 identificado como primer recorte (ADR-011) & No (en seguimiento) \\
    R-02 & Fallos de seguridad en el código de partida &
      Auditoría adversarial de RLS con pruebas pgTAP y refutación independiente; pruebas de regresión (ADR-015) & Sí: B-05--B-09 \\
    R-03 & Exposición del entorno de desarrollo (servidor remoto con IP pública) &
      Filtrado en la cadena \texttt{DOCKER-USER}, acceso por túnel SSH y revisión de cada servicio nuevo (ADR-010) & Sí: B-04 \\
    R-04 & Errores introducidos por cambios masivos o por agentes de IA &
      Controles automáticos ampliados a las variantes del error, revisión adversarial y prueba manual del autor & Sí: B-01, B-14, B-16 \\
    R-05 & Deriva entre el esquema declarativo y las migraciones &
      \texttt{supabase db diff} en la verificación de todo cambio de BD & Sí: B-10 \\
    R-06 & Pruebas que no prueban lo que dicen &
      Leer las aserciones de las pruebas heredadas, no solo sus nombres & Sí: B-11 \\
    R-07 & Interrupción de las herramientas o de la conexión durante una tarea &
      Cerrar cada incremento con su \textit{commit}; tareas de seguridad formuladas como verificación defensiva y divididas en etapas & Sí: B-12, B-15 \\
    R-08 & Falsos positivos de las herramientas automáticas (secretos, pruebas inestables) &
      Documentar las credenciales públicas de desarrollo; ejecutar los E2E contra la \textit{build} de test & Sí: B-03, B-13 \\
    R-09 & Conflicto de licencia del código de partida &
      Repositorio privado y origen documentado internamente (ADR-001, ADR-006) & No \\
    R-10 & Validación de la reutilización sin un dominio de negocio concreto &
      Medir el arranque de un SaaS nuevo a partir de la plataforma (ADR-005, punto abierto P-01) & Pendiente \\
    \bottomrule
  \end{tabularx}
\end{table}

\todo{Asignar probabilidad e impacto (p.\,ej.\ bajo/medio/alto) a cada riesgo
según el criterio del autor y, si la plantilla lo pide, una matriz de
riesgos.}

Entre los riesgos materializados destacan tres por su enseñanza para la
planificación:

\begin{itemize}
  \item \textbf{B-04, exposición del entorno.} El entorno de desarrollo es un
        servidor remoto. Los servicios locales de la base de datos publicaban
        sus puertos en todas las interfaces y el motor de contenedores
        insertaba reglas que se saltaban el cortafuegos, de modo que la API y
        la base de datos (con credenciales de desarrollo conocidas) fueron
        accesibles desde Internet durante aproximadamente una hora, solo con
        datos de prueba. Se corrigió con una regla persistente, se recrearon
        los volúmenes y se pasó a acceder mediante un túnel SSH (ADR-010).
        Además, dio lugar a un criterio de revisión para cualquier servicio
        nuevo y a una advertencia que figurará en el manual de instalación.
  \item \textbf{B-05 a B-09, fallos heredados.} La auditoría de la base de
        datos del CMS encontró dos brechas confirmadas y varios
        endurecimientos necesarios en el código de partida (se analizan en la
        sección~\ref{sec:verificacion-seguridad}). El riesgo se había previsto
        y la mitigación (auditoría obligatoria en todo cambio de BD) fue la
        que permitió detectarlos.
  \item \textbf{B-15, corte de conexión con un agente trabajando.} El trabajo
        quedó a medias sin \textit{commit} y mezclado con una corrección de
        seguridad pendiente. La lección se incorporó al proceso: cada
        incremento se cierra con su propio \textit{commit}.
\end{itemize}

% ---------------------------------------------------------------------
\section{Herramientas}
\label{sec:herramientas}

\begin{itemize}
  \item \textbf{Control de versiones: git y GitHub.} Repositorio privado
        (condición de la licencia del código de partida, ADR-001). La rama
        principal es \texttt{main} y se trabaja en ramas por fase
        (\texttt{fase-1/base}, \texttt{fase-2/cms}). Los mensajes de
        \textit{commit} están en español y siguen la especificación
        \textit{Conventional Commits}~\cite{conventionalcommits}, con el
        módulo como ámbito.
  \item \textbf{Integración continua: GitHub Actions.} Un único
        \textit{workflow} con tres trabajos: (1)~controles propios del TFG
        (el \textit{scope} de paquetes es bloqueante; el desmarcado completo
        y el informe de comentarios en inglés son informativos hasta cerrar
        F3 y F4); (2)~comprobación de tipos, \textit{lint}, formato y pruebas
        unitarias; y (3)~base de datos local, pruebas pgTAP y E2E con
        Playwright contra la \textit{build} de test. Para ahorrar minutos de
        ejecución, el tercero solo se lanza en las \textit{pull requests}
        hacia \texttt{main} y a mano.
  \item \textbf{Monorepo:} Turborepo~\cite{turborepo} y pnpm, con un catálogo
        único de versiones.
  \item \textbf{Base de datos local:} Supabase CLI~\cite{supabase-docs} sobre
        Docker.
  \item \textbf{Pruebas:} Vitest (unitarias), pgTAP~\cite{pgtap} (base de
        datos) y Playwright~\cite{playwright} (E2E).
  \item \textbf{Asistente de programación:} Claude Code~\cite{claudecode},
        con el \textit{harness} descrito en la
        sección~\ref{sec:metodologia-ia}.
  \item \textbf{Scripts propios de control} (\texttt{scripts/tfg/}):
        comprobación de desmarcado, informe de comentarios en inglés y
        extracción de las etiquetas de trazabilidad del código.
  \item \textbf{Memoria:} \LaTeX{} sobre la plantilla de la ESI en Overleaf.
        % PENDIENTE (P-04): confirmar al integrar la plantilla oficial.
\end{itemize}

% ---------------------------------------------------------------------
\section{Costes}
\label{sec:costes}

% PENDIENTE (autor): todos los importes. No hay datos de costes en el
% repositorio; esta sección solo propone la estructura.

\todo{Completar la estimación de costes con datos reales.}

\begin{table}[htbp]
  \centering
  \small
  \caption{Estructura de costes del proyecto (pendiente de completar).}
  \label{tab:costes}
  \begin{tabularx}{\textwidth}{@{}Xrr@{}}
    \toprule
    \textbf{Concepto} & \textbf{Base de cálculo} & \textbf{Importe (€)} \\
    \midrule
    Personal: horas de desarrollo $\times$ coste/hora de un ingeniero junior & \textit{pendiente: horas y tarifa} & \textit{pendiente} \\
    Equipo de desarrollo (amortización) & \textit{pendiente} & \textit{pendiente} \\
    Servidor remoto de desarrollo (VPS) & \textit{pendiente: meses} & \textit{pendiente} \\
    Licencia comercial de la base de código de partida % PENDIENTE (P-06): cómo se presenta
      & \textit{pendiente} & \textit{pendiente} \\
    Suscripción al asistente de programación con IA & \textit{pendiente: meses} & \textit{pendiente} \\
    Servicios en la nube para el despliegue (F7) & \textit{pendiente} & \textit{pendiente} \\
    Otros (dominio, etc.) & \textit{pendiente} & \textit{pendiente} \\
    \midrule
    \textbf{Total} & & \textit{pendiente} \\
    \bottomrule
  \end{tabularx}
\end{table}

Las herramientas de desarrollo empleadas (git, Node.js, pnpm, Turborepo,
Supabase CLI, Docker, Vitest, pgTAP, Playwright) son de código abierto y no
tienen coste de licencia. Stripe se usa en modo de pruebas.
\todo{Confirmar el plan contratado de GitHub, del proyecto de Supabase en la
nube y de la plataforma de despliegue, y si generan coste.}

\setcounter{chapter}{3}
% =====================================================================
% Capítulo 4 · Análisis
% Borrador generado a partir de docs/tfg/REQUISITOS.md, TRAZABILIDAD.md,
% DECISIONES.md (ADR-004, ADR-005, ADR-011, ADR-014) y AGENTS.md.
% Estado de las fuentes: 2026-09-29 (F2.4b cerrada).
% Paquetes necesarios: booktabs, tabularx, todonotes, tikz
%   (bibliotecas: positioning, shapes.geometric, fit, backgrounds).
% Alternativa en PlantUML: diagramas/casos-uso.puml
% =====================================================================

\chapter{Análisis}
\label{cap:analisis}

Este capítulo delimita el problema que resuelve PymeKit, recoge sus
requisitos funcionales y no funcionales y describe los casos de uso con los
que se validará la plataforma.

% ---------------------------------------------------------------------
\section{Contexto del problema}
\label{sec:contexto-problema}

Las pequeñas y medianas empresas (pymes) consumen cada vez más aplicaciones en
modalidad \textit{Software as a Service} (SaaS, software como servicio): el
proveedor aloja y mantiene una única instalación y cada cliente accede a ella
mediante una suscripción. Quien desarrolla un SaaS para pymes debe resolver,
antes de escribir una sola línea de su lógica de negocio, un conjunto de
problemas que se repiten en casi todos los productos de este tipo:

\begin{itemize}
  \item \textbf{identidad}: registro, inicio de sesión, recuperación de
        contraseña y, cada vez más a menudo, un segundo factor de
        autenticación;
  \item \textbf{multi-tenencia} (\textit{multi-tenancy}): cada pyme cliente
        es un \textit{tenant} (inquilino) cuyos datos deben quedar aislados de
        los de los demás, aunque todos compartan la misma base de datos;
  \item \textbf{organización interna del cliente}: equipos, invitaciones,
        roles y permisos;
  \item \textbf{cobro}: planes, suscripciones y sincronización con una
        pasarela de pago;
  \item \textbf{operación de la plataforma}: una consola para que el
        personal del proveedor gestione usuarios y cuentas, y una herramienta
        para consultar y corregir datos sin acceder directamente a la base de
        datos;
  \item \textbf{calidad transversal}: seguridad, internacionalización,
        pruebas y despliegue reproducible.
\end{itemize}

Estos componentes no aportan valor diferencial al producto, pero consumen una
parte muy importante del esfuerzo y son precisamente los que más riesgo de
seguridad concentran. El objetivo de PymeKit es ofrecerlos ya resueltos,
integrados y verificados, como una \textbf{plataforma base reutilizable}: un
nuevo SaaS para pymes debería poder arrancarse cambiando la configuración, sin
modificar el núcleo (RNF-01).

Del alcance acordado se derivan dos decisiones que condicionan el análisis:

\begin{itemize}
  \item La plataforma es \textbf{genérica}: no implementa un dominio de
        negocio concreto. Su reutilización se valida midiendo el arranque de
        un SaaS nuevo a partir de ella (ADR-005).
        % PENDIENTE (P-01): confirmar con el tutor que esto cubre el
        % «escenario de ejemplo» de la propuesta.
  \item El \textbf{CMS} que pide la propuesta se interpreta como un CMS de
        administración de datos de la base de datos, integrado en la consola
        de administración, y no como un gestor de contenidos editoriales
        (blog, documentación). Mantener dos CMS añadiría complejidad sin
        acercar el objetivo (ADR-004, ADR-011).
\end{itemize}

% PENDIENTE (P-06): si la memoria explica aquí que la plataforma parte de una
% base de código SaaS existente con licencia, redactarlo de forma neutra (sin
% nombres comerciales ni URL) según lo que se acuerde con el tutor. La parte
% reutilizada y la propia se distinguen en el capítulo de Implementación.

% ---------------------------------------------------------------------
\section{Actores}
\label{sec:actores}

\begin{description}
  \item[Visitante.] Persona no identificada que consulta la web pública y
        puede registrarse.
  \item[Usuario.] Persona registrada. Dispone de una cuenta personal y puede
        pertenecer a una o varias cuentas de equipo.
  \item[Propietario del equipo.] Usuario con el rol de mayor nivel en una
        cuenta de equipo (la pyme cliente); gestiona sus miembros y su
        suscripción.
  \item[Miembro.] Usuario de un equipo con un rol de menor nivel y, por
        tanto, con permisos limitados.
  \item[Super-administrador.] Personal del proveedor con control total de la
        plataforma y, por extensión, del CMS (ADR-014).
  \item[Gestor de contenidos / soporte.] Personal del proveedor que accede al
        CMS con permisos limitados definidos por el RBAC propio del CMS
        (ADR-014, ADR-016).
  \item[Desarrollador.] Quien reutiliza PymeKit para construir un SaaS nuevo.
  \item[Pasarela de pago.] Sistema externo (Stripe) que procesa los cobros y
        notifica los cambios de suscripción mediante \textit{webhooks}
        (llamadas HTTP que el sistema externo envía a la aplicación cuando
        ocurre un evento).
\end{description}

% ---------------------------------------------------------------------
\section{Requisitos funcionales}
\label{sec:rf}

Los requisitos se derivan de los apartados \emph{Objetivo}, \emph{Descripción}
y \emph{Alcance} de la propuesta del TFG. Su prioridad es \textbf{M}
(imprescindible, forma parte del prototipo mínimo) o \textbf{D} (deseable).
Los identificadores son estables: se citan en el código, en la matriz de
trazabilidad y en esta memoria, y no se renumeran.

\begin{table}[htbp]
  \centering
  \footnotesize
  \caption{Requisitos funcionales.}
  \label{tab:rf}
  \begin{tabularx}{\textwidth}{@{}lXc@{}}
    \toprule
    \textbf{ID} & \textbf{Requisito} & \textbf{Prior.} \\
    \midrule
    RF-01 & Interfaz web para el usuario final: página de inicio pública (\textit{landing}), página de precios y área privada con navegación & M \\
    RF-02 & API de \textit{backend} para la lógica de negocio y el acceso a datos, mediante funciones de servidor (\textit{server functions}) tipadas y validadas & M \\
    RF-03 & Registro e inicio de sesión (contraseña y enlace mágico), recuperación de contraseña y verificación del correo electrónico & M \\
    RF-04 & Segundo factor de autenticación (MFA) y proveedores OAuth configurables & D \\
    RF-05 & Gestión de la cuenta personal: perfil, correo, contraseña, preferencia de idioma y baja & M \\
    RF-06 & Cuentas de equipo (una por pyme cliente): crear equipos, invitar miembros, asignar roles y aplicar permisos por funcionalidad & M \\
    RF-07 & Suscripciones con pasarela de pago (Stripe): planes, pago (\textit{checkout}), portal de cliente y sincronización mediante \textit{webhooks} & M \\
    RF-08 & Panel de super-administración de la plataforma: listar y gestionar usuarios y cuentas, bloquear, suplantar y restablecer & M \\
    RF-09 & CMS de administración de datos, integrado en la consola de administración: explorar y editar tablas de la BD, gestionar usuarios y almacenamiento, con permisos propios & M \\
    RF-10 & Registro de auditoría de las acciones realizadas desde el CMS & D \\
    RF-11 & Paneles (\textit{dashboards}) configurables en el CMS & D \\
    RF-12 & Notificaciones dentro de la aplicación y correos transaccionales & D \\
    RF-13 & Interfaz multidioma: español por defecto e inglés & M \\
    \bottomrule
  \end{tabularx}
\end{table}

RF-09 es el requisito que más ha evolucionado durante el análisis: en un
principio se planteó el CMS como una aplicación separada que compartía la
base de datos (ADR-002) y, tras evaluar alternativas, se decidió integrarlo
en la consola de administración (ADR-011). De ahí que su redacción incluya
«integrado en la consola de administración». RF-11 es el primer candidato a
recortar si el plazo lo exige.

% ---------------------------------------------------------------------
\section{Requisitos no funcionales}
\label{sec:rnf}

Cada requisito no funcional se acompaña del modo en que se verifica, para que
sea comprobable y no una mera declaración de intenciones.

\begin{table}[htbp]
  \centering
  \footnotesize
  \caption{Requisitos no funcionales y forma de verificarlos.}
  \label{tab:rnf}
  \begin{tabularx}{\textwidth}{@{}l>{\raggedright\arraybackslash}p{6.2cm}X@{}}
    \toprule
    \textbf{ID} & \textbf{Requisito} & \textbf{Verificación} \\
    \midrule
    RNF-01 & \textbf{Reutilización y configurabilidad}: la plataforma se adapta a un nuevo SaaS cambiando la configuración (ficheros de configuración y variables de entorno), sin tocar el núcleo & Evaluación de la reutilización (F6): pasos y tiempo para arrancar un proyecto nuevo \\
    RNF-02 & \textbf{Seguridad}: la autorización se aplica en la base de datos (RLS) además de en la aplicación; secretos fuera del repositorio y validación de todas las entradas & Pruebas pgTAP y revisiones adversariales de RLS y del código \\
    RNF-03 & \textbf{Aislamiento multi-tenant}: los datos de una cuenta no son accesibles desde otra & Pruebas pgTAP específicas de aislamiento \\
    RNF-04 & \textbf{Modularidad}: monorepo de paquetes desacoplados, con dependencias explícitas y exportaciones separadas para cliente y servidor & Estructura del monorepo, comprobación de dependencias y de tipos \\
    RNF-05 & \textbf{Mantenibilidad}: TypeScript estricto, \textit{lint} y formato automáticos, y código comentado en español didáctico & Integración continua y revisión de comentarios \\
    RNF-06 & \textbf{Calidad verificable}: pruebas unitarias, de base de datos y E2E de los casos de uso principales & Integración continua en verde \\
    RNF-07 & \textbf{Desplegabilidad}: despliegue reproducible con contenedores y en la nube, con manual de instalación & Despliegue de prueba (F7) \\
    RNF-08 & \textbf{Usabilidad y accesibilidad}: interfaz adaptable (\textit{responsive}), componentes accesibles y tema claro/oscuro & Revisión manual y E2E \\
    RNF-09 & \textbf{Rendimiento razonable}: carga de datos en el servidor, caché de consultas y paginación en el CMS & Revisión manual \\
    \bottomrule
  \end{tabularx}
\end{table}

La relación entre cada requisito, el código que lo implementa y las pruebas
que lo verifican se mantiene en una matriz de trazabilidad que se actualiza al
cerrar cada incremento.
\todo{Incluir la matriz de trazabilidad (requisito $\rightarrow$ módulo
$\rightarrow$ pruebas $\rightarrow$ estado) como anexo o en el capítulo de
Pruebas, con el estado final tras F6.}

% ---------------------------------------------------------------------
\section{Casos de uso}
\label{sec:casos-uso}

Se han seleccionado siete casos de uso principales que recorren los módulos
de la plataforma y que servirán de base para la validación de la fase F6
(tabla~\ref{tab:cu} y figura~\ref{fig:casos-uso}).

\begin{table}[htbp]
  \centering
  \footnotesize
  \caption{Casos de uso principales.}
  \label{tab:cu}
  \begin{tabularx}{\textwidth}{@{}l>{\raggedright\arraybackslash}p{2.6cm}Xl@{}}
    \toprule
    \textbf{CU} & \textbf{Actor} & \textbf{Descripción} & \textbf{Requisitos} \\
    \midrule
    CU-01 & Visitante & Se registra, verifica el correo y accede al área privada & RF-01, RF-03 \\
    CU-02 & Usuario & Crea un equipo para su pyme e invita a un empleado con el rol \textit{member} & RF-06 \\
    CU-03 & Propietario del equipo & Contrata un plan (pasarela en modo de pruebas) y el estado de la suscripción se refleja en la aplicación & RF-07 \\
    CU-04 & Miembro sin permiso & Intenta gestionar la facturación y se le deniega, tanto en la interfaz como en la base de datos & RF-06, RNF-02, RNF-03 \\
    CU-05 & Super-administrador & Busca un usuario, lo bloquea y lo suplanta para darle soporte & RF-08 \\
    CU-06 & Gestor de contenidos & Entra en el CMS y edita registros de una tabla según sus permisos; la acción queda auditada & RF-09, RF-10 \\
    CU-07 & Desarrollador & Arranca un nuevo SaaS a partir de PymeKit siguiendo el manual & RNF-01, RNF-07 \\
    \bottomrule
  \end{tabularx}
\end{table}

\begin{figure}[htbp]
  \centering
  \begin{tikzpicture}[
      font=\footnotesize,
      uc/.style={ellipse, draw, minimum width=4.2cm, minimum height=0.85cm,
                 align=center, inner sep=1pt},
      actor/.style={draw=none, align=center, font=\footnotesize\bfseries},
      link/.style={-}
    ]
    % Casos de uso (dentro del sistema)
    \node[uc] (cu1) at (0, 3.6)  {CU-01 Registrarse y\\acceder al área privada};
    \node[uc] (cu2) at (0, 2.4)  {CU-02 Crear equipo\\e invitar a un miembro};
    \node[uc] (cu3) at (0, 1.2)  {CU-03 Contratar\\un plan};
    \node[uc] (cu4) at (0, 0)    {CU-04 Gestionar la\\facturación (denegado)};
    \node[uc] (cu5) at (0, -1.2) {CU-05 Bloquear y\\suplantar a un usuario};
    \node[uc] (cu6) at (0, -2.4) {CU-06 Editar registros\\en el CMS (auditado)};
    \node[uc] (cu7) at (0, -3.6) {CU-07 Arrancar un\\SaaS nuevo};

    \begin{scope}[on background layer]
      \node[draw, rounded corners, fit=(cu1)(cu7), inner sep=10pt,
            label={[font=\small\bfseries]above:PymeKit}] (sys) {};
    \end{scope}

    % Actores (izquierda: clientes; derecha: personal y externos)
    \node[actor] (vis) at (-5.2, 3.6)  {Visitante};
    \node[actor] (usu) at (-5.2, 2.4)  {Usuario};
    \node[actor] (pro) at (-5.2, 1.2)  {Propietario\\del equipo};
    \node[actor] (mie) at (-5.2, 0)    {Miembro\\sin permiso};
    \node[actor] (sa)  at (5.2, -1.2)  {Super-\\administrador};
    \node[actor] (ges) at (5.2, -2.4)  {Gestor de\\contenidos};
    \node[actor] (dev) at (5.2, -3.6)  {Desarrollador};
    \node[actor] (pas) at (5.2, 1.2)   {Pasarela\\de pago};

    \draw[link] (vis) -- (cu1);
    \draw[link] (usu) -- (cu2);
    \draw[link] (pro) -- (cu3);
    \draw[link] (mie) -- (cu4);
    \draw[link] (sa)  -- (cu5);
    \draw[link] (ges) -- (cu6);
    \draw[link] (dev) -- (cu7);
    \draw[link] (pas) -- (cu3);
  \end{tikzpicture}
  \caption{Diagrama de casos de uso principales de PymeKit.}
  \label{fig:casos-uso}
  % Texto alternativo: diagrama con el sistema PymeKit y siete casos de uso
  % (CU-01 a CU-07). A la izquierda, los actores cliente (visitante, usuario,
  % propietario del equipo, miembro sin permiso) enlazados con CU-01 a CU-04;
  % a la derecha, el personal (super-administrador con CU-05, gestor de
  % contenidos con CU-06), el desarrollador con CU-07 y la pasarela de pago
  % como actor secundario de CU-03.
\end{figure}
% PENDIENTE: si la plantilla prefiere figuras de actor UML (monigote), usar la
% versión PlantUML de diagramas/casos-uso.puml y exportarla a PDF/SVG.

A continuación se describe brevemente cada caso de uso.

\begin{description}
  \item[CU-01 · Registro y acceso.] El visitante se registra con correo y
        contraseña (o con enlace mágico), recibe el correo de verificación,
        lo confirma y accede al área privada, donde dispone de su cuenta
        personal. Cubre RF-01 y RF-03.
  \item[CU-02 · Equipo e invitación.] Un usuario crea una cuenta de equipo
        para su pyme, de la que pasa a ser propietario, e invita por correo a
        un empleado con el rol \textit{member}. El invitado acepta la
        invitación y queda como miembro con los permisos de su rol. Cubre
        RF-06.
  \item[CU-03 · Contratación de un plan.] El propietario elige un plan y
        completa el pago en la pasarela (en modo de pruebas). La pasarela
        notifica el resultado mediante un \textit{webhook} y la aplicación
        refleja el estado de la suscripción de la cuenta. Cubre RF-07.
  \item[CU-04 · Acceso denegado.] Un miembro sin el permiso de gestión de la
        facturación intenta acceder a ella. La interfaz no se lo ofrece y,
        aunque forzara la petición, la base de datos la rechaza por sus
        políticas de seguridad. Es el caso que materializa la seguridad en
        profundidad (RNF-02) y el aislamiento (RNF-03).
  \item[CU-05 · Soporte desde la consola.] El super-administrador, con sesión
        de doble factor, busca un usuario en la consola, lo bloquea y lo
        suplanta temporalmente para reproducir un problema. Cubre RF-08.
  \item[CU-06 · Edición en el CMS.] Un gestor de contenidos, con un rol
        limitado del CMS, abre una tabla para la que tiene permiso, edita un
        registro y la acción queda registrada en la auditoría. Las tablas sin
        permiso no le aparecen ni puede consultarlas. Cubre RF-09 y RF-10.
  \item[CU-07 · Reutilización.] Un desarrollador arranca un SaaS nuevo a
        partir de PymeKit siguiendo el manual y cambiando solo la
        configuración. Sirve para evaluar RNF-01 y RNF-07.
        % PENDIENTE (P-01): definir el escenario y las métricas concretas.
\end{description}

\todo{Cuando termine F6, añadir para cada caso de uso su flujo principal y
alternativo (plantilla de caso de uso de la ESI, si la hay) y la prueba que lo
valida. En la fecha de este borrador, CU-01 está probado a mano por el autor;
CU-02 a CU-06 dependen de F2 y F5, y CU-07 de F6.}

% =====================================================================
% Capítulo 5 · Diseño (borrador PARCIAL: solo lo decidido e implementado)
% Fuentes: código de apps/web, packages/cms y apps/web/supabase/schemas;
% docs/tfg/DECISIONES.md (ADR-001…016), BITACORA.md (B-04…B-09, B-14),
% TRAZABILIDAD.md, PROGRESO.md; packages/cms/AGENTS.md;
% .claude/skills/rls-review.
% Estado de las fuentes: 2026-09-29 (F2.4b cerrada).
% Paquetes necesarios: booktabs, tabularx, todonotes, tikz
%   (bibliotecas: positioning, arrows.meta, fit, backgrounds, calc).
% Alternativas en PlantUML: diagramas/arquitectura.puml,
%   diagramas/flujo-peticion-cms.puml
% =====================================================================

\chapter{Diseño}
\label{cap:diseno}

% PENDIENTE: borrador parcial. Faltan el modelo de datos completo, pagos
% (RF-07, F5), roles orientados a pymes (F5), configurabilidad (F5) y el
% diseño de las pantallas del CMS que aún no están hechas (F2.4c–F2.8).

Este capítulo describe el diseño de PymeKit en lo que ya está decidido e
implementado: la arquitectura general, el modelo multi-tenant y de control de
acceso, la estrategia de seguridad en profundidad y el proceso con el que se
ha verificado esa seguridad. Cada decisión remite al ADR en que se justifica.

% PENDIENTE (P-06): este capítulo presenta el diseño resultante. La plataforma
% se construyó partiendo de una base de código SaaS existente, con licencia, y
% de un CMS para el mismo tipo de base de datos (ADR-001). Acordar con el
% tutor cómo se indica aquí qué partes del diseño son heredadas y cuáles
% propias; como orientación (docs/tfg/MAPA-REFERENCIAS.md): el modelo de
% cuentas y el RBAC de equipos y del CMS son heredados; la integración del CMS
% en la web (ADR-011), el pegamento super-admin/CMS (ADR-014), el acceso a la
% consola (ADR-016) y las correcciones de seguridad (ADR-015) son propios.

% ---------------------------------------------------------------------
\section{Arquitectura general}
\label{sec:arquitectura}

\subsection{Monorepo de paquetes}

El código se organiza en un único \textit{monorepo} (un repositorio con
varias aplicaciones y paquetes que se versionan juntos) gestionado con
Turborepo~\cite{turborepo} y pnpm (ADR-002). Contiene:

\begin{itemize}
  \item \texttt{apps/web}: la aplicación web, que incluye tanto la
        aplicación SaaS como la consola de administración y el CMS;
  \item \texttt{apps/web/supabase}: los esquemas, migraciones, datos de
        prueba y pruebas pgTAP de \emph{toda} la base de datos;
  \item \texttt{apps/e2e}: las pruebas E2E con Playwright;
  \item \texttt{packages/*}: los paquetes reutilizables de la plataforma
        (\texttt{@pymekit/*}: interfaz, autenticación, cuentas, equipos,
        facturación, administración, i18n\ldots) y, en
        \texttt{packages/cms/*}, los del CMS (\texttt{@pymekit/cms-*}).
\end{itemize}

En la fecha de este borrador el monorepo tiene 47 paquetes, sin ciclos de
dependencias (RNF-04). Los paquetes declaran sus dependencias de forma
explícita y separan las exportaciones de cliente y de servidor: por ejemplo,
los paquetes del CMS con rutas y servicios son solo de servidor, y el código
de la interfaz del CMS vive en paquetes cliente aparte
(\texttt{@pymekit/cms-ui-core}, \texttt{@pymekit/cms-data-explorer-ui}\ldots).
Esta separación no es solo una cuestión de orden: evita que código con
credenciales de servidor acabe en el navegador y, como se comprobó en la
incidencia B-16, también evita ciclos entre paquetes. La \textit{build} de la
web incluye además un control que falla si el cliente de base de datos del
servidor aparece en un fragmento (\textit{chunk}) destinado al navegador.
% PENDIENTE: el nº de paquetes (47) procede de TRAZABILIDAD (RNF-04); volver a
% contarlo al cerrar F2.

\subsection{Pila tecnológica}

\begin{itemize}
  \item \textbf{Aplicación web:} TanStack Start~\cite{tanstack-start} con
        React 19. Proporciona enrutado por ficheros (TanStack
        Router~\cite{tanstack-router}), renderizado en el servidor
        (\textit{Server-Side Rendering}, SSR), cargadores de datos por ruta
        (\textit{loaders}) y \textit{server functions}: funciones que se
        ejecutan en el servidor y se invocan desde el cliente con tipos
        compartidos (RF-02). Las cachés de datos del cliente usan TanStack
        Query~\cite{tanstack-query}.
  \item \textbf{Datos, autenticación y almacenamiento:}
        Supabase~\cite{supabase-docs}, que ofrece PostgreSQL, un servicio de
        autenticación basado en JWT y almacenamiento de ficheros.
  \item \textbf{API del CMS:} Hono~\cite{hono} (\textit{framework} HTTP
        ligero basado en el estándar \texttt{Request}/\texttt{Response}) y
        Drizzle ORM~\cite{drizzle} para el acceso a la base de datos.
  \item \textbf{Pagos:} Stripe (RF-07). % PENDIENTE: diseño de pagos en F5.
  \item \textbf{Interfaz:} Tailwind CSS 4 y componentes accesibles propios
        del paquete \texttt{@pymekit/ui} (RNF-08).
\end{itemize}

Para que el conjunto se mantenga coherente se decidió usar una sola pila de
librerías en toda la web, CMS incluido: un único sistema de formularios, uno
de internacionalización y una sola biblioteca de componentes (ADR-013). Así
cada cosa se hace de una única forma (RNF-05) y el idioma español solo hay que
añadirlo en un sistema (RF-13).

\subsection{Una única base de datos}

La aplicación y el CMS comparten \textbf{una sola instancia} de PostgreSQL
(ADR-002). Los datos de la plataforma viven en el esquema \texttt{public} y
los del CMS (su RBAC, sus preferencias, las vistas guardadas, la auditoría y
los paneles) en un esquema propio llamado \texttt{cms} (ADR-012). Compartir la
base de datos es lo que permite que el CMS administre los datos de la propia
plataforma (RF-09) y que las mismas políticas de seguridad protejan ambas
partes.

\subsection{CMS integrado en la consola de administración}

La decisión de mayor calado del diseño es cómo integrar el CMS (ADR-011). Se
evaluaron tres opciones:

\begin{enumerate}[label=(\alph*)]
  \item \textbf{aplicaciones separadas} (la web, la interfaz del CMS y su
        API como tres servicios): la de menor esfuerzo, pero con tres
        despliegues, dos inicios de sesión y las \textit{cookies} de sesión
        duplicadas;
  \item \textbf{integración híbrida} (la API dentro de la web y la interfaz
        del CMS servida como aplicación estática): convivirían dos
        enrutadores y dos formas de programar;
  \item \textbf{integración total}: la API se monta dentro del servidor de
        la web y las pantallas se reescriben como rutas de la web.
\end{enumerate}

Se eligió la opción (c), a pesar de ser la de mayor esfuerzo estimado, porque
da un único servicio, un único inicio de sesión y un único estilo de código.
En concreto:

\begin{itemize}
  \item La aplicación Hono del CMS se monta en \texttt{/api/cms/*} mediante
        una ruta de servidor de TanStack Start (\texttt{routes/api/cms/\$.ts})
        que le reenvía la petición tal cual. El CMS comparte así proceso,
        dominio y \textit{cookies} con la web: un super-administrador que ya
        ha iniciado sesión en la consola lo usa sin volver a identificarse.
  \item Las pantallas del CMS son rutas de la web bajo
        \texttt{/admin/cms/*}, con \textit{loaders} y TanStack Query sobre un
        cliente RPC tipado de la API.
  \item Durante el SSR, los \textit{loaders} no hacen una petición de red a
        la propia API: llaman a la aplicación Hono dentro del mismo proceso,
        reenviando solo la cabecera \texttt{cookie} de la petición entrante
        (ADR-016).
\end{itemize}

La figura~\ref{fig:arquitectura} resume la arquitectura resultante.

\begin{figure}[htbp]
  \centering
  \begin{tikzpicture}[
      font=\footnotesize,
      box/.style={draw, rounded corners, align=center, minimum height=0.9cm,
                  inner sep=4pt},
      group/.style={draw, dashed, rounded corners, inner sep=8pt},
      arr/.style={-{Stealth[length=2mm]}}
    ]
    % Cliente
    \node[box, minimum width=10cm] (browser) at (0, 5.2)
      {Navegador (React 19): área privada, consola \texttt{/admin} y CMS \texttt{/admin/cms}};

    % Servidor web
    \node[box, minimum width=3.0cm] (ssr) at (-3.6, 3.3)
      {Rutas y \textit{loaders}\\(SSR, TanStack Router)};
    \node[box, minimum width=3.0cm] (sfn) at (0, 3.3)
      {\textit{Server functions}\\+ \textit{middleware}};
    \node[box, minimum width=3.0cm] (hono) at (3.6, 3.3)
      {Ruta \texttt{/api/cms/*}\\$\rightarrow$ aplicación Hono};

    \node[box, minimum width=3.0cm] (sbclient) at (-1.8, 1.7)
      {Cliente Supabase\\(sesión del usuario)};
    \node[box, minimum width=3.0cm] (drizzle) at (3.6, 1.7)
      {Servicios Drizzle\\(\textit{claims} del usuario)};

    \begin{scope}[on background layer]
      \node[group, fit=(ssr)(sfn)(hono)(sbclient)(drizzle),
            label={[font=\small\bfseries]left:\rotatebox{90}{apps/web (TanStack Start)}}] (web) {};
    \end{scope}

    % Supabase
    \node[box, minimum width=2.6cm] (auth) at (-3.6, -0.4) {Auth\\(JWT, MFA)};
    \node[box, minimum width=4.4cm] (pg) at (0.4, -0.4)
      {PostgreSQL + RLS\\esquemas \texttt{public} y \texttt{cms}};
    \node[box, minimum width=2.2cm] (st) at (4.2, -0.4) {Storage};
    \begin{scope}[on background layer]
      \node[group, fit=(auth)(pg)(st),
            label={[font=\small\bfseries]below:Supabase (una única instancia)}] (sb) {};
    \end{scope}

    % Externo
    \node[box] (stripe) at (-7.6, 1.7) {Stripe};

    \draw[arr] (browser) -- (ssr);
    \draw[arr] (browser) -- (sfn);
    \draw[arr] (browser) -- node[right, font=\scriptsize]{RPC} (hono);
    \draw[arr] (ssr) -- (sfn);
    \draw[arr, dashed] (ssr.east) ++(0,-0.25) -- node[below, font=\scriptsize]{SSR en proceso} ([yshift=-0.25cm]hono.west);
    \draw[arr] (sfn) -- (sbclient);
    \draw[arr] (hono) -- (drizzle);
    \draw[arr] (sbclient) -- (pg);
    \draw[arr] (sbclient) -- (auth);
    \draw[arr] (drizzle) -- (pg);
    \draw[arr] (stripe.east) -- node[above, font=\scriptsize]{\textit{webhooks}} (stripe.east -| web.west);
  \end{tikzpicture}
  \caption{Arquitectura general de PymeKit: una única aplicación web con el
    CMS integrado y una única instancia de Supabase.}
  \label{fig:arquitectura}
  % Texto alternativo: el navegador se comunica con la aplicación web
  % (TanStack Start), que contiene las rutas con SSR, las server functions y
  % la ruta /api/cms que monta la aplicación Hono del CMS. Las server
  % functions usan el cliente Supabase con la sesión del usuario; la API del
  % CMS usa servicios Drizzle con los claims del usuario. Ambos acceden a una
  % única instancia de Supabase (Auth, PostgreSQL con RLS en los esquemas
  % public y cms, y Storage). Stripe notifica a la web mediante webhooks.
\end{figure}
% PENDIENTE: revisar el ajuste fino de posiciones al compilar (no se ha
% podido compilar en el entorno donde se generó el borrador). Versión
% alternativa en diagramas/arquitectura.puml.
% Nota: los webhooks de Stripe entran por la ruta de servidor
% /api/billing/webhook de la web (apps/web/src/routes/api/billing/webhook.ts);
% el diseño de pagos se completa en F5.

% ---------------------------------------------------------------------
\section{Modelo multi-tenant y control de acceso}
\label{sec:multitenant-rbac}

\subsection{Cuentas personales y de equipo}

El \textit{tenant} de PymeKit es la \textbf{cuenta}. La tabla
\texttt{accounts} contiene dos clases de cuenta, distinguidas por el campo
\texttt{is\_personal\_account}:

\begin{itemize}
  \item la \textbf{cuenta personal}, una por usuario, cuyo identificador
        coincide con el del usuario en el servicio de autenticación;
  \item la \textbf{cuenta de equipo}, pensada para representar a una pyme
        cliente (RF-06), con un propietario principal, miembros
        (\texttt{accounts\_memberships}) e invitaciones.
\end{itemize}

Los datos de negocio se enlazan a una cuenta mediante una columna
\texttt{account\_id}, y las políticas de seguridad a nivel de fila deciden el
acceso en función de la relación del usuario con esa cuenta. Cada miembro
tiene un rol con un nivel jerárquico, y cada rol tiene un conjunto de
permisos por funcionalidad (\texttt{roles.manage}, \texttt{billing.manage},
\texttt{settings.manage}, \texttt{members.manage}, \texttt{invites.manage}).
Las políticas se apoyan en funciones auxiliares como
\texttt{has\_role\_on\_account} (pertenencia) y \texttt{has\_permission}
(permiso concreto). Las escrituras destructivas se condicionan al permiso, no
a la mera pertenencia.
\todo{Añadir el diagrama entidad-relación del esquema \texttt{public} y, tras
F5, los roles orientados a pymes.}

\subsection{Control de acceso del CMS}

El CMS tiene su propio control de acceso basado en roles
(\textit{Role-Based Access Control}, RBAC) en el esquema \texttt{cms}: cuentas
del CMS, roles con rango, grupos de permisos y permisos por tabla o por
almacenamiento. Conviven así dos modelos de autorización en la misma base de
datos: el de la plataforma (el super-administrador y los equipos) y el del
CMS. La decisión de diseño (ADR-014) fue conectarlos en lugar de sustituir
uno por otro:

\begin{itemize}
  \item \textbf{El super-administrador es la raíz del CMS.} Unos
        \textit{triggers} sobre los usuarios mantienen sincronizados a los
        super-administradores con un rol de sistema \emph{Root}, el de mayor
        rango, que concede todos los permisos. Al promocionar a un usuario
        recibe el acceso al CMS, una cuenta activa y ese rol; al retirarle la
        condición pierde el acceso, y su cuenta del CMS se desactiva pero se
        conserva para la auditoría. Un índice único garantiza que solo exista
        un rol raíz.
  \item \textbf{El RBAC del CMS da acceso limitado a otro personal}
        (soporte, gestión de contenidos). Se descartó dejar solo al
        super-administrador porque incumpliría RF-09 («con permisos propios»).
\end{itemize}

Ambos grupos comparten la consola de administración, pero con límites
distintos (ADR-016): la consola admite a quien es super-administrador
\emph{o} tiene acceso al CMS; las páginas de gestión de la plataforma exigen
super-administrador, y el personal del CMS que intente abrirlas se redirige
al CMS. La barra lateral solo muestra las secciones que la API permite, pero
esa ocultación es solo de interfaz: cada ruta y cada \textit{endpoint} vuelven
a comprobar el permiso.

% ---------------------------------------------------------------------
\section{Seguridad en profundidad}
\label{sec:seguridad}

El principio que guía el diseño de la seguridad es que \textbf{la base de
datos es la frontera de autorización} (RNF-02): toda comprobación en la
interfaz o en el servidor de aplicación se considera una primera barrera, útil
para dar una buena experiencia, pero nunca suficiente. Si una capa superior
falla, las políticas de la base de datos siguen impidiendo el acceso. Este
enfoque de defensa en profundidad es coherente con las recomendaciones de
verificación de seguridad de aplicaciones de OWASP~\cite{owasp-asvs}.
% VERIFICAR: si se cita ASVS, referenciar los requisitos concretos (control
% de acceso, V4 en la versión 4.0.3) tras comprobar la versión vigente.

\subsection{RLS como frontera}

PostgreSQL permite asociar a cada tabla políticas de seguridad a nivel de
fila (\textit{Row-Level Security}, RLS)~\cite{postgresql-rls}: predicados que
el propio motor añade a cada consulta según el usuario que la ejecuta.
Supabase~\cite{supabase-rls} expone la base de datos a clientes
autenticados y traslada la identidad del usuario (su JWT) a la sesión de
PostgreSQL, de modo que las políticas pueden decidir con ella. En PymeKit:

\begin{itemize}
  \item toda tabla tiene RLS activado; se revocan todos los privilegios por
        defecto a los roles de la API y se conceden solo los verbos que la
        funcionalidad usa (y el \texttt{UPDATE}, por columnas, para que no se
        pueda mover una fila a otra cuenta);
  \item las políticas de \texttt{UPDATE} repiten en su cláusula
        \texttt{WITH CHECK} las comprobaciones del \texttt{INSERT}, porque sin
        ella PostgreSQL solo valida la fila antigua~\cite{postgresql-createpolicy}
        (véase B-05);
  \item las funciones \texttt{security definer}, que se ejecutan con
        privilegios elevados y se saltan RLS, deben aplicar sus propias
        comprobaciones de acceso y justificarlas en un comentario;
  \item el cliente de servidor que ignora RLS se reserva para casos
        excepcionales, tras validar a mano los permisos.
\end{itemize}

\subsection{MFA obligatorio en la consola y en el CMS}

El servicio de autenticación distingue el nivel de garantía de la sesión
(\textit{Authenticator Assurance Level}, AAL)~\cite{supabase-mfa}:
\texttt{aal1} tras la contraseña y \texttt{aal2} tras completar el segundo
factor. El diseño exige \texttt{aal2} en todas las capas que dan acceso a la
administración:

\begin{itemize}
  \item en la base de datos, la función que identifica al super-administrador
        devuelve falso si la sesión no es \texttt{aal2}; unas políticas
        restrictivas exigen \texttt{aal2} a quien tiene un factor configurado;
        y el acceso al CMS lo decide la función
        \texttt{cms.verify\_admin\_access()}, que comprueba la cuenta activa
        del CMS y el MFA;
  \item en la interfaz, la consola \texttt{/admin} exige siempre
        \texttt{aal2}: con un factor configurado redirige a la verificación y
        vuelve después a la página pedida; sin factor responde 404 (ADR-016).
\end{itemize}

En PymeKit el MFA es \textbf{obligatorio por defecto} en el CMS: una opción
de configuración ausente o no válida se interpreta como «obligatorio».

\subsection{Middleware más comprobación en la base de datos}

La figura~\ref{fig:flujo-cms} muestra el recorrido de una petición a la API
del CMS y las comprobaciones que encuentra en cada capa:

\begin{enumerate}
  \item \textbf{Guardas de ruta} de la consola (\texttt{/admin} y
        \texttt{/admin/cms}): sesión \texttt{aal2} y acceso de administración.
  \item \textbf{Ruta \texttt{/api/cms/\$}} de TanStack Start: reenvía la
        petición a Hono. El código del CMS se importa dinámicamente dentro
        del manejador para que nunca llegue al navegador.
  \item \textbf{Hono}: cabeceras de seguridad, protección CSRF propia (la
        global de la web solo cubre las \textit{server functions}) y
        \textit{middleware} de autenticación, que verifica el JWT (401 sin
        sesión), exige el \textit{claim} de acceso al CMS (403) y que el
        usuario no esté bloqueado.
  \item \textbf{Servicio Drizzle}: cada transacción fija los
        \textit{claims} del usuario y el rol \texttt{authenticated} en la
        sesión de PostgreSQL mediante parámetros enlazados (nunca se
        interpola un valor del JWT en el SQL), de modo que las políticas RLS
        del esquema \texttt{cms} se aplican igual que a cualquier cliente.
  \item \textbf{PostgreSQL}: \texttt{cms.verify\_admin\_access()} y las
        políticas RLS toman la decisión final. Un super-administrador con
        sesión \texttt{aal1} supera el \textit{middleware}, pero no ve datos.
\end{enumerate}

\begin{figure}[htbp]
  \centering
  \begin{tikzpicture}[
      font=\footnotesize,
      step/.style={draw, rounded corners, align=center, minimum width=3.3cm,
                   minimum height=1.0cm, inner sep=3pt},
      check/.style={align=left, font=\scriptsize, text width=6.4cm},
      arr/.style={-{Stealth[length=2mm]}}
    ]
    \node[step] (s1) at (0, 0)     {Navegador\\(consola \texttt{/admin/cms})};
    \node[step] (s2) at (0, -1.6)  {Ruta de servidor\\\texttt{/api/cms/\$}};
    \node[step] (s3) at (0, -3.2)  {Aplicación Hono\\(\textit{middleware})};
    \node[step] (s4) at (0, -4.8)  {Servicio Drizzle\\(transacción)};
    \node[step] (s5) at (0, -6.4)  {PostgreSQL\\(esquema \texttt{cms})};

    \draw[arr] (s1) -- (s2);
    \draw[arr] (s2) -- (s3);
    \draw[arr] (s3) -- (s4);
    \draw[arr] (s4) -- (s5);

    \node[check, right=0.6cm of s1] {Guarda de ruta: sesión \texttt{aal2} y
      acceso de administración (si no: verificación MFA o 404)};
    \node[check, right=0.6cm of s2] {Reenvío a Hono; código del CMS importado
      solo en el servidor};
    \node[check, right=0.6cm of s3] {Cabeceras de seguridad, CSRF; JWT válido
      (401), \textit{claim} de acceso al CMS (403), usuario no bloqueado};
    \node[check, right=0.6cm of s4] {\textit{Claims} del usuario y rol
      \texttt{authenticated} fijados con parámetros enlazados};
    \node[check, right=0.6cm of s5] {\texttt{cms.verify\_admin\_access()}
      (cuenta activa, MFA) y políticas RLS: \textbf{decisión final}};
  \end{tikzpicture}
  \caption{Flujo de una petición a la API del CMS y comprobaciones de
    seguridad en cada capa.}
  \label{fig:flujo-cms}
  % Texto alternativo: diagrama vertical de cinco pasos (navegador, ruta de
  % servidor /api/cms/$, aplicación Hono, servicio Drizzle, PostgreSQL). A la
  % derecha de cada paso se indican sus comprobaciones; la última, las
  % políticas RLS y cms.verify_admin_access(), es la que toma la decisión.
\end{figure}
% Versión alternativa como diagrama de secuencia: diagramas/flujo-peticion-cms.puml

Las \textit{server functions} de la aplicación siguen el mismo patrón: una
cadena de \textit{middleware} tipado (autenticación, pertenencia al equipo,
super-administrador) valida la petición y después el cliente de Supabase con
la sesión del usuario deja que RLS decida.

\subsection{Fallo en cerrado}

Todas las comprobaciones de seguridad deben \textbf{fallar en cerrado}
(\textit{fail closed}): ante un dato ausente, ilegible o no válido, se
deniega. El diseño lo aplica, por ejemplo, a la opción que indica si el CMS
exige MFA, que se lee con una función de privilegios elevados precisamente
para que RLS no pueda ocultársela a quien tiene que decidir. Esta regla se
incorporó tras la incidencia B-06 (sección~\ref{sec:verificacion-seguridad}).

\subsection{Seguridad del entorno de desarrollo}

La seguridad también alcanza al entorno de trabajo. Los servicios locales de
la base de datos se aíslan de Internet con una regla de cortafuegos
específica para contenedores y se accede a la aplicación mediante un túnel
SSH (ADR-010, tras la incidencia B-04). Cualquier servicio nuevo en
contenedores se revisa con ese mismo criterio.

% ---------------------------------------------------------------------
\section{Decisiones de diseño principales}
\label{sec:decisiones}

La tabla~\ref{tab:decisiones} resume las decisiones de diseño que se han
tomado hasta la fecha, con las alternativas descartadas.

\begin{table}[htbp]
  \centering
  \footnotesize
  \caption{Resumen de decisiones de diseño (ADR) y alternativas descartadas.}
  \label{tab:decisiones}
  \begin{tabularx}{\textwidth}{@{}l>{\raggedright\arraybackslash}p{4.3cm}>{\raggedright\arraybackslash}Xl@{}}
    \toprule
    \textbf{ADR} & \textbf{Decisión} & \textbf{Alternativas descartadas (motivo)} & \textbf{Req.} \\
    \midrule
    001 & Partir de una base de código SaaS existente con licencia % PENDIENTE (P-06)
      & Desarrollo propio completo (inviable en plazo); otros \textit{boilerplates} abiertos (menos completos en multi-tenant y RLS); CMS genéricos como Strapi o Directus (otra base de datos u otro \textit{backend}) & RNF-01, 04 \\
    002 & Monorepo único y una sola base de datos & Dos repositorios con la misma BD (CI, despliegue y documentación duplicados; peor trazabilidad) & RF-09, RNF-04 \\
    004 & Conservar el CMS de datos y retirar el de contenidos editoriales & Mantener dos CMS (complejidad sin aportar al objetivo) & RF-09 \\
    010 & Aislar de Internet los servicios locales & Confiar en el cortafuegos del sistema (los contenedores se lo saltan, B-04) & RNF-02 \\
    011 & CMS integrado en la consola: API Hono en \texttt{/api/cms}, pantallas como rutas de la web & Apps separadas (tres servicios, dos inicios de sesión); híbrida (dos enrutadores) & RF-08--11 \\
    012 & Esquema SQL del CMS llamado \texttt{cms} & Conservar el nombre original (excepción permanente en el control de desmarcado) & RF-09 \\
    013 & Una sola pila de librerías en toda la web & Mantener las del CMS (dos formas de hacer lo mismo; dos sistemas de i18n) & RNF-05 \\
    014 & Super-admin como raíz del CMS; RBAC del CMS para el resto del personal & Solo super-admin, sin RBAC granular (incumple RF-09) & RF-08, 09 \\
    015 & Endurecer la seguridad del CMS de partida & Aceptar el código heredado sin auditar (brechas confirmadas, B-05--B-09) & RNF-02, 03 \\
    016 & Consola para super-admin y personal del CMS, siempre con \texttt{aal2} & Dejar entrar con el \textit{claim} de acceso sin comprobar el AAL (B-14) & RF-08, RNF-02 \\
    \bottomrule
  \end{tabularx}
\end{table}

% ---------------------------------------------------------------------
\section{Verificación de la seguridad}
\label{sec:verificacion-seguridad}

\subsection{Proceso de revisión adversarial}

Todo cambio que afecta a la base de datos (una política, un permiso, una
función \texttt{security definer}, una vista o una migración) se somete a una
revisión adversarial específica de RLS, definida como un procedimiento del
\textit{harness} (sección~\ref{sec:metodologia-ia}). Su premisa es que leer
las políticas no basta: hay que \textbf{demostrar} el aislamiento. Consta de
estas etapas:

\begin{enumerate}
  \setcounter{enumi}{-1}
  \item \textbf{Alcance}: qué tablas, funciones y esquemas se revisan.
  \item \textbf{Inventario de aislamiento}: políticas, permisos concedidos,
        funciones de privilegios elevados, vistas y restricciones de MFA y
        super-administrador de cada objeto.
  \item \textbf{Auditoría estática} con un catálogo de ataques (exposición de
        tablas, huecos en los predicados, capas restrictivas, funciones que
        se saltan RLS, vistas, almacenamiento y el esquema del CMS). Cada
        acierto se formula como un ataque concreto: qué usuario, qué
        sentencia y qué fila.
  \item \textbf{Demostración con pgTAP}~\cite{pgtap}: para cada hallazgo
        plausible, y para cada garantía que se quiere confirmar, se escribe
        una prueba que ejecuta el ataque \emph{como el usuario atacante} y
        comprueba que falla. Las pruebas quedan en el repositorio como
        pruebas de regresión.
  \item \textbf{Refutación independiente}: el contexto que produjo un
        hallazgo no puede confirmarlo. Un agente distinto recibe solo la
        política, el resultado de la prueba y la afirmación, y su única tarea
        es refutarla (encontrar la capa que en realidad protege o, ante un
        veredicto de «seguro», el tipo de sentencia, rol o nivel de sesión
        que no se ha probado). Una brecha solo se acepta si el refutador
        también la reproduce; un aprobado, solo si no encuentra un vector
        abierto.
  \item \textbf{Informe} con el veredicto, los hallazgos, la matriz de
        aislamiento y las correcciones.
\end{enumerate}

\subsection{Hallazgos corregidos}

La aplicación de este proceso a la base de datos del CMS en el incremento
F2.1 dio los resultados de la tabla~\ref{tab:hallazgos}, que también recoge
dos incidencias de seguridad detectadas por otras vías. Los fallos de origen
\emph{heredado} ya estaban en el código de partida; los \emph{propios} se
introdujeron durante el proyecto.

\begin{table}[htbp]
  \centering
  \footnotesize
  \caption{Incidencias de seguridad detectadas y corregidas (bitácora del proyecto).}
  \label{tab:hallazgos}
  \begin{tabularx}{\textwidth}{@{}l>{\raggedright\arraybackslash}Xll>{\raggedright\arraybackslash}p{3.1cm}@{}}
    \toprule
    \textbf{ID} & \textbf{Fallo} & \textbf{Origen} & \textbf{Grav.} & \textbf{Detección} \\
    \midrule
    B-04 & Base de datos local accesible desde Internet (los contenedores se saltaban el cortafuegos) & entorno & Alta & Investigación de un problema de acceso \\
    B-05 & Escalada de privilegios en el RBAC del CMS: \texttt{UPDATE} sin \texttt{WITH CHECK} & heredado & Alta & Auditoría, confirmada con prueba ejecutada \\
    B-06 & La comprobación de MFA del CMS fallaba en abierto & heredado & Alta & Una prueba nueva falló de forma inesperada \\
    B-07 & Lectura de la tabla de usuarios de autenticación mediante una función de consulta & heredado & Alta & Auditoría \\
    B-08 & Lectura del catálogo de roles de PostgreSQL (\texttt{pg\_authid}) & heredado & Alta & Refutación independiente \\
    B-09 & Política RLS con una condición tautológica & heredado & Baja & Auditoría \\
    B-14 & La consola \texttt{/admin} se cargaba sin segundo factor & propio & Media & Prueba manual del autor \\
    \bottomrule
  \end{tabularx}
\end{table}
% PENDIENTE: la bitácora no indica expresamente cómo se detectaron B-07 y
% B-09 (se asume «auditoría», dentro de la /rls-review de F2.1). Confirmar.

A continuación se resume cada uno.

\begin{description}
  \item[B-05 · Escalada por \texttt{UPDATE} sin \texttt{WITH CHECK}.] Las
        políticas de actualización de tres tablas de asignación de permisos
        solo tenían la cláusula \texttt{USING}. Un administrador delegado
        podía cambiar el permiso de una fila existente y colgar de un rol
        inferior un permiso que él mismo no tenía. Se añadió \texttt{WITH
        CHECK} con la misma comprobación que el \texttt{INSERT}. Lección: en
        PostgreSQL, un \texttt{UPDATE} sin \texttt{WITH CHECK} valida la fila
        antigua, no la nueva.
  \item[B-06 · MFA que fallaba en abierto.] La función de control leía la
        opción «MFA obligatorio» con los permisos del usuario, y una política
        restrictiva se la ocultaba precisamente a quien tenía el factor
        configurado pero había entrado sin él. La función veía la opción
        vacía, la interpretaba como «opcional» y dejaba pasar. Se corrigió
        leyendo la opción con una función de privilegios elevados y tratando
        su ausencia como «obligatorio». Al corregirlo falló una prueba
        heredada cuyo nombre decía «MFA obligatorio por defecto» pero cuya
        aserción comprobaba que se concedía el acceso (B-11).
  \item[B-07 · Lectura de usuarios de autenticación.] Dos funciones de
        consulta que se saltan RLS no aplicaban la validación de esquemas
        protegidos ni, una de ellas, la de acceso de administración, que sí
        aplicaban sus funciones hermanas de escritura. Con un permiso
        comodín podían devolver los usuarios con sus \textit{hashes} de
        contraseña. Ahora exigen ambas comprobaciones.
  \item[B-08 · Lectura del catálogo del sistema.] Tras corregir B-07, la
        etapa de refutación encontró que la lista de esquemas protegidos no
        incluía el catálogo de PostgreSQL, lo que permitía leer los
        \textit{hashes} de los roles de la base de datos. Se bloquea ahora
        cualquier esquema de sistema. Es la evidencia más clara del valor de
        la refutación independiente: encontró algo que la primera auditoría
        había dado por bueno.
  \item[B-09 · Condición tautológica.] En una subconsulta dentro de una
        política, una columna sin cualificar se resolvía contra la tabla de
        la subconsulta y se comparaba consigo misma, de modo que cualquier
        usuario con un rol veía todas las asignaciones. Se cualificó la
        columna.
  \item[B-14 · Consola sin segundo factor.] Al abrir la consola al personal
        del CMS, la guarda comprobaba el \textit{claim} de acceso, que
        también tiene un super-administrador que solo ha introducido la
        contraseña. La API y la base de datos seguían exigiendo MFA, así que
        no se exponían datos, pero la consola se cargaba. Lo detectó el autor
        con una prueba exploratoria; se corrigió exigiendo \texttt{aal2} en
        la guarda y se añadió una prueba E2E de regresión. Es también un
        ejemplo de la defensa en profundidad funcionando: el fallo de una
        capa no expuso datos porque las inferiores lo impidieron.
\end{description}

Las correcciones de la base de datos se recogen en ADR-015 y en tres
migraciones, cada una con sus pruebas de regresión en pgTAP. ADR-015 enumera
ocho puntos: dos brechas confirmadas (B-05 y B-08), cinco endurecimientos
(el fallo en abierto de la comprobación de MFA, B-06; la lectura de esquemas
protegidos, B-07; el acceso a los paneles; la unicidad de los marcadores del
rol raíz, y la política tautológica, B-09) y la alineación de cuatro funciones
cuyo esquema declarativo no coincidía con las migraciones (B-10).
Tras la corrección, la batería de pruebas de base de datos quedó en 57
ficheros y 1.550 pruebas pgTAP en verde.
% Fuente: PROGRESO.md (F2.1). Actualizar al cierre de F2/F6.

El mismo ADR deja registradas, para revisarlas en F2.7, cuatro debilidades
menores que hoy no constituyen una brecha activa, como que un permiso de
almacenamiento sin contenedor indicado equivalga a un comodín o que el
personal pueda insertar entradas de auditoría a su nombre.
\todo{Actualizar este párrafo con el resultado de la revisión en F2.7.}

\todo{En el capítulo de Pruebas, presentar la matriz de aislamiento por tabla
del informe de la revisión y las cifras finales de pruebas (pgTAP, unitarias
y E2E) tras F6.}


\bibliographystyle{plain}
\bibliography{referencias}
\end{document}
